\subsection{Challenge for spontaneous CP violation in Type IIB orientifolds with fluxes --- arXiv:2004.04518}

\textbf{Authors:} Tatsuo Kobayashi, Hajime Otsuka

\paragraph{主な主張}
解析したType IIBトーラスorientifoldでは、3形式フラックスによるCP保存真空とCP破れ真空の間に平坦方向が残り、フラックスだけではCPを破る孤立真空の選択に不十分であることを示した。 \cite{Kobayashi:2020uaj}

\paragraph{新規性}
CP不変なフラックス配置を系統的に調べ、CP破れの障害を真空の縮退として特定するとともに、特別な配置では4次元CPがモジュラー双対性へ埋め込まれることを示した。

\paragraph{位置付け}
フラックス安定化だけで自発的CP破れを実現できる範囲を明確にし、平坦方向を持ち上げる追加効果の検討を動機付ける研究である。
